\documentclass[11pt,a4paper]{article}
\usepackage[margin=2.4cm]{geometry}
\usepackage{amsmath,amssymb,bm}
\usepackage{booktabs}
\usepackage{graphicx}
\usepackage{hyperref}
\usepackage{xcolor}
\usepackage{microtype}

\newcommand{\ee}{E_\theta}
\newcommand{\dd}{\mathrm{d}}
\newcommand{\code}[1]{\texttt{#1}}

\title{The SnapWave wave model coupled to SFINCS:\\[2pt]
\large Governing equations as implemented in the source code,\\
and comparison with Boussinesq and spectral wave models}
\author{Technical note based on a source-code review of\\
\texttt{SFINCS/source/src/snapwave/} and \texttt{sfincs\_snapwave.f90}\\
(checkout \texttt{cline/f61jy65y}, SFINCS v2.4.0 Galibier+)}
\date{\today}

\begin{document}
\maketitle

\begin{abstract}
\noindent
This note documents the wave-model equations actually implemented in SnapWave,
the stationary wave solver embedded in SFINCS, as read directly from the Fortran
source code. We derive the governing balance equations, list every source and
sink term with its parameters and code defaults, describe the numerical scheme
and the two-way coupling with the SFINCS hydrodynamics, and then compare the
formulation with phase-resolving Boussinesq-type models and full spectral wave
models (SWAN/XBeach-style action balance). The goal is to provide a precise
reference for diagnosing unexpected behaviour in coupled SFINCS--SnapWave
simulations.
\end{abstract}

\tableofcontents
\newpage

\section{Model context and overall structure}

SnapWave is a \emph{stationary}, \emph{phase-averaged}, \emph{directionally
resolved} wave model. It does not resolve individual wave phases or the full
frequency spectrum. Instead, at each node $k$ of an unstructured mesh it
carries a one-dimensional directional distribution of wave energy density
$\ee(\theta,k)$ (units J\,m$^{-2}$\,rad$^{-1}$) at a single representative
frequency $\sigma(k) = 2\pi/T_p(k)$ per wave system. Two wave systems are
solved separately:

\begin{itemize}
  \item \textbf{Incident (short) waves}: energy $\ee$, height $H$, period $T_p$.
  \item \textbf{Infragravity (IG) waves}: energy $E_{\theta,ig}$, height $H_{ig}$,
  period $T_{p,ig}$, with their own wavenumber, group velocity, refraction speed
  and dissipation terms (suffix \code{\_ig} in the code).
\end{itemize}

The solver is called from SFINCS every \code{dtwave} seconds (default 3600\,s)
through \code{sfincs\_snapwave.f90::update\_wave\_field}. Before each call,
SFINCS water levels $\zeta$ are mapped onto the SnapWave mesh so that the wave
model uses the instantaneous depth $h = \max(\zeta - z_b,\,10^{-5}\,\mathrm{m})$.
After the call, wave heights, periods and depth-integrated wave forces are
mapped back onto the SFINCS grid and added to the momentum equation. SnapWave
itself is stationary: time enters only through the boundary conditions and the
slowly varying water depth, and a pseudo-time term is used purely as an
iteration device (Section~\ref{sec:numerics}).

\section{Governing equation: the stationary directional energy balance}
\label{sec:governing}

\subsection{Continuous form}

The core equation, implemented in
\code{snapwave\_solver.f90::solve\_energy\_balance2Dstat}, is a stationary,
directionally resolved wave-energy balance at representative frequency
$\sigma$:

\begin{equation}
\underbrace{\frac{\partial \ee}{\partial t^{*}}}_{\text{pseudo-time}}
+\underbrace{c_g\,\frac{\partial \ee}{\partial s}}_{\text{propagation}}
+\underbrace{\frac{\partial}{\partial \theta}\!\left(c_\theta\,\ee\right)}_{\text{refraction}}
=
-\underbrace{\frac{D_{w}+D_{f}+D_{veg}}{E}\,\ee}_{\text{dissipation}}
+\underbrace{S_{\mathrm{wind}}}_{\text{optional}}
\;\mp\;
\underbrace{S_{ig}}_{\substack{\text{incident--IG}\\\text{exchange}}},
\label{eq:energybalance}
\end{equation}

where $\ee(\theta)$ is the directionally resolved wave-energy density
(J\,m$^{-2}$\,rad$^{-1}$), $s$ is the along-ray coordinate, $c_g$ the group
velocity, $c_\theta$ the refraction (direction-turning) speed, and
$D_w,D_f,D_{veg}$ the breaking, bottom-friction and vegetation dissipation
rates (W\,m$^{-2}$). The total energy is
\begin{equation}
E=\int_{\theta}\ee\,\dd\theta,\qquad
H=\sqrt{\frac{8E}{\rho g}} \;=\; H_{\mathrm{rms}},
\label{eq:Hrms}
\end{equation}
so the $H$ carried internally by SnapWave is the \emph{root-mean-square} wave
height. When coupling back to SFINCS and to the boundary input, the code
converts with $H_{m0}=\sqrt{2}\,H_{\mathrm{rms}}$
(\code{sfincs\_snapwave.f90}: \code{hm0 = snapwave\_H * sqrt(2.0)}).

Equation~\eqref{eq:energybalance} is the energy form of the wave-action
balance: because frequency is \emph{not} a coordinate and no current field is
applied in SnapWave, action conservation reduces to energy conservation and the
balance is written in energy form directly. This is a fundamental simplification
relative to SWAN-type models (Section~\ref{sec:comparison}).

An identical equation with IG coefficients ($c_{g,ig}$, $c_{\theta,ig}$,
$D_{w,ig}$, $D_{f,ig}$) is solved for $E_{\theta,ig}$, with the sign of the
exchange term reversed: the term acts as a \emph{sink} for incident energy and a
\emph{source} for IG energy (scaled by \code{snapwave\_shinc2ig}, default 1.0).

\subsection{Linear wave kinematics}

The dispersion relation is evaluated with Eckart's (1952) explicit
approximation (\code{snapwave\_windsource.f90::disper\_approx\_1}):
\begin{equation}
k=\frac{\sigma^{2}}{g}
\left[1-\exp\!\left(-\left(\sigma\sqrt{h/g}\right)^{2.5}\right)\right]^{-0.4},
\qquad
c=\frac{\sigma}{k},\qquad
n=\tfrac{1}{2}+\frac{kh}{\sinh(2kh)},\qquad
c_g=n\,c .
\label{eq:dispersion}
\end{equation}
The refraction speed follows from linear ray theory for a depth-varying,
current-free medium:
\begin{equation}
c_\theta(\theta)=
\frac{\sigma}{\sinh(2kh)}
\left(\frac{\partial h}{\partial x}\sin\theta-\frac{\partial h}{\partial y}\cos\theta\right),
\qquad
|c_\theta|\le \frac{\sigma}{4},
\label{eq:ctheta}
\end{equation}
where the limiter (one quarter turn per wave period) prevents unrealistic
turning rates over steep bathymetry. Note that $\partial h/\partial x$ is
evaluated from the \emph{bed} gradient field computed once from $z_b$ at
initialisation (\code{snapwave\_domain.f90::fm\_surrounding\_points}), while
$h$ itself varies in time with the SFINCS water level. There is \textbf{no
refraction by currents}: ambient currents $\mathbf{u}$ from SFINCS are not
passed to the ray equations.


\section{Boundary conditions and directional spectra}
\label{sec:boundaries}

Offshore input consists of bulk parameters $H_{m0}$, $T_p$, mean direction
$\bar\theta_w$ (nautical, coming-from, degrees) and directional spreading
$\sigma_\theta$ (degrees) at one or more support points
(\code{snapwave\_boundaries.f90}). At each support point the total energy is set
from the significant height,
\begin{equation}
E_0=\frac{1}{16}\rho g H_{m0}^{2},
\end{equation}
and distributed over directions with a wrapped power-of-cosine distribution
centred on the mean wave direction:
\begin{equation}
\ee(\theta)\propto
\begin{cases}
\operatorname{sgn}(\cos(\theta-\bar\theta))\,
|\cos(\theta-\bar\theta)|^{m_s},
& |\theta-\bar\theta|\lesssim \pi/2,\\[2pt]
0, & \text{otherwise},
\end{cases}
\qquad m_s=\frac{1}{\sigma_\theta^{2}}-1 .
\label{eq:cosn}
\end{equation}
Grid boundary cells receive spectra by inverse-distance weighting of the two
nearest support points. On lateral (Neumann) boundaries all wave properties are
copied from the adjacent inner cell. The directional grid $\{\theta_j\}$ is
rebuilt \emph{at every SnapWave call} around the current energy-weighted mean
wave (or wind) direction (\code{make\_theta\_grid}), so the directional window
moves in time; the sector width is \code{snapwave\_sector} (default
$180^\circ$) with resolution \code{snapwave\_dtheta} (default $10^\circ$).

\section{Dissipation and source terms}
\label{sec:sources}

\subsection{Depth-induced breaking (Baldock et al., 1998)}

Breaking dissipation (\code{snapwave\_solver.f90::baldock}) is
\begin{equation}
D_w=\frac{0.28\,\alpha\,\rho g}{T}\,
\exp\!\left[-\left(\frac{H_{\max}}{H}\right)^{2}\right]
\left(H_{\max}^{2}+H^{2}\right)\,f ,
\qquad
f=\left(\frac{H}{H_{\max}}\right)^{i_{\exp}}\ \text{if } H>H_{\max},
\label{eq:baldock}
\end{equation}
with $H_{\max}=\gamma h$ for incident waves ($\gamma$ =
\code{snapwave\_gamma}, default 0.7; $\alpha$ = \code{snapwave\_alpha},
default 1.0), and the Battjes--Janssen-type limiter
$H_{\max,ig}=(0.88/k_{ig})\tanh(\gamma_{ig} k_{ig} h/0.88)$ for IG waves. The
extra factor $f$ (exponent \code{snapwave\_baldock\_exponent}, default 2) adds
dissipation on steep coasts where Baldock's form cannot keep up with shoaling.
Breaking is only evaluated when $H>$~\code{snapwave\_baldock\_ratio}~$\times
H_{\max}$ (default ratio 0.2).

\subsection{Bottom friction}

For incident waves, a quadratic friction law with the near-bed orbital velocity
amplitude gives
\begin{equation}
D_f=0.28\,\rho\, f_w\, u_{orb}^{3},
\qquad
u_{orb}=\frac{\sigma H}{2\sinh(kh)},
\end{equation}
with $f_w$ = \code{snapwave\_fw} (default 0.01). For IG waves the code uses a
formulation proportional to the incident wave height
(\code{Df\_ig} in \code{solve\_energy\_balance2Dstat}):
\begin{equation}
D_{f,ig}=0.0361\,f_{w,ig}\left(\frac{g}{h}\right)^{3/2} H_{\mathrm{inc}}\,E_{ig},
\qquad f_{w,ig}=\code{snapwave\_fwig}\ (0.015).
\end{equation}

\subsection{Vegetation}

With \code{vegetation=1}, wave dissipation by vegetation $D_{veg}$ follows the
layered Suzuki et al.\ (2012) formulation with a Keulegan--Carpenter-number bulk
drag coefficient (adapted from XBeach; \code{vegatt}, \code{swvegatt},
\code{bulkdragcoeff}), using vertically segmented stem height $a_h$, diameter
$b_v$, density $N$ and drag $C_D$ from \code{vegetationfile}. An additional
direct wave-drag force on the mean flow, $F_{vw}$, is computed from a
nonlinear (skewed/asymmetric) wave-shape velocity model (\code{swvegnonlin},
\code{momeqveg}).

\subsection{Depth limiter}

After every node solve, energy exceeding the depth limit is removed by uniform
directional scaling:
\begin{equation}
\ee \leftarrow \ee \Big/
\max\!\left[1,\;
\left(\frac{\sqrt{E/\left(\tfrac18\rho g\right)}}{\gamma_{\max} h}\right)^{2}
\right],
\qquad \gamma_{\max}=\code{snapwave\_gammax}\ (\text{default } 999).
\label{eq:depthlimit}
\end{equation}
This is a hard limiter distinct from the Baldock term; with the default
$\gamma_{\max}=999$ it is effectively inactive unless set by the user.


\section{Infragravity wave generation and propagation}
\label{sec:ig}

\subsection{Boundary condition: Herbers et al.\ (1994) bound-wave spectrum}

If \code{snapwave\_use\_herbers=1} (default), the offshore IG boundary
condition is computed from the incident sea state: a JONSWAP frequency spectrum
(peak enhancement $\gamma_J$ = \code{snapwave\_jonswapgamma}, default 3.3) is
built from $H_{m0}$, $T_p$ and spreading, and the bound long-wave energy is
computed from the bispectrum-based closure of Herbers et al.\ (1994), following
the XBeach implementation (\code{snapwave\_infragravity.f90::compute\_herbers}).
The IG period at the boundary is chosen from the reconstructed IG spectrum by
\code{snapwave\_tpig\_opt} (1=$T_{m01}$, default). Stability guards in
\code{determine\_ig\_bc}: if $H_{m0}/h>0.5$ at the support point the depth is
reset to $2H_{m0}$; depths are capped at 200\,m. With
\code{snapwave\_use\_herbers=0}, the older parameterisation
$E_{ig,0}=\code{snapwave\_eeinc2ig}\times E_{inc}$ (default ratio 0.01) and
$T_{ig}= \code{snapwave\_Tinc2ig}\times T_p$ (default 7) is used instead.

\subsection{Inshore IG source: release of the bound long wave}

Along the propagation path, IG energy is generated from the shoaling of the
incident-wave radiation stress following Leijnse et al.\ (2024)
(\code{determine\_infragravity\_source\_sink\_term}). Per direction the source
inserted into the IG right-hand side is
\begin{equation}
S_{ig}(\theta)=
\alpha_{ig}\,
\frac{\sqrt{E_{ig}^{\,\uparrow}}\; c_g^{\,\uparrow}}{h^{\uparrow}}\,
\frac{\max(\Delta S_{xx},0)}{\Delta s}\,
\frac{\ee(\theta,k)}{E(k)},
\label{eq:srcig}
\end{equation}
where $\uparrow$ denotes upwind-interpolated values. For
\code{snapwave\_ig\_opt=1} (default), $\Delta S_{xx}$ uses a conservative
shoaling estimate,
\begin{equation}
S_{xx}^{\mathrm{cons}}=
E^{\uparrow}\,\frac{c_g^{\,\uparrow}}{c_{g,ig}(k)}\,
\bigl(2n(k)-\tfrac12\bigr),
\end{equation}
relative to its upwind value; \code{ig\_opt=2} takes the raw difference. The
efficiency factor $\alpha_{ig}$ is an empirical function of local bed slope
$\beta$ and $\gamma_{\mathrm{loc}}=H_{\mathrm{rms}}/h$
(\code{estimate\_shoaling\_parameter\_alphaig}, fit constants as in the code),
limited to $[0,1]$ and scaled by \code{snapwave\_alphaigfac} (default 1.0). The
source is smoothly switched off once incident waves break significantly,
$\gamma_{\mathrm{loc}} >$ \code{snapwave\_gamma\_fac\_br}~$\times\,\gamma$
(default 0.45), via logistic transition factors; beyond that point IG waves are
treated as free. The same fraction \code{snapwave\_shinc2ig} of the source is
subtracted from the incident balance (default: all of it).

\section{Wind input and variable wave period}
\label{sec:wind}

With \code{snapwave\_wind=1}, a coupled \emph{wave-action} balance
$A=\ee/\sigma$ is solved alongside Eq.~\eqref{eq:energybalance}, allowing the
representative period to evolve: after each node solve,
$\sigma=\operatorname{minmax}(E/A,\,\sigma_{\min},\,\sigma_{\max})$
(\code{snapwave\_sigmin}, \code{snapwave\_sigmax} bound $\sigma$ to
$[2\pi/25,\,2\pi/1]$\,rad\,s$^{-1}$ by default). Wind input follows the
dimensionless growth curves of Breugem and Holthuijzen (2007) in
Kahma--Calkoen form, capped by fully developed (Pierson--Moskowitz) and
depth-limited sea states, and distributed over directions with
\code{windspreadfac} (\code{snapwave\_windsource.f90::windinput}).
Dissipation in the action equation uses the same fractional rate as the energy
equation, with a $T$-dependent correction (\code{DwT}, \code{DwAk} terms).


\section{Wave forces and the two-way coupling with SFINCS}
\label{sec:coupling}

\subsection{Dissipation-based wave force}

SnapWave does \emph{not} evaluate the full radiation-stress tensor gradient
$-\partial S_{ij}/\partial x_j$. Instead, the wave force per unit mass applied
to the water column is computed from the local dissipation
(\code{solve\_energy\_balance2Dstat}, final loop):
\begin{equation}
F(k)=\underbrace{\frac{D_w(k)\,k}{\sigma\,\rho\, h}}_{\approx\, D_w/(\rho h c)}
\;+\;\underbrace{\frac{D_{veg}(k)\,k}{\sigma\,\rho\, h}}_{\text{vegetation dissipation}}
\;+\;\underbrace{F_{vw}(k)}_{\substack{\text{direct wave drag}\\\text{from vegetation}}},
\qquad [\mathrm{m\,s^{-2}}],
\label{eq:force}
\end{equation}
i.e.\ the force is approximated as dissipation divided by (density $\times$
depth $\times$ phase speed), directed along the local energy-weighted mean
direction $\bar\theta(k)=\operatorname{atan2}(\sum\ee\sin\theta,\sum\ee\cos\theta)$.
This is the usual ``$D/c$'' closure used in stationary surf-zone models; it
assumes that the divergence of the radiation-stress field is locally balanced
by dissipation, which is a reasonable approximation in breaking-dominated surf
zones but not in shoaling or reflecting regions where $\nabla\!\cdot\!\mathbf{F}_{rad}$
is set by wave--wave interaction and geometry rather than by $D$.

\subsection{Transfer to the momentum equation}

In \code{sfincs\_snapwave.f90::compute\_snapwave}, the force is converted to a
depth-integrated form $\mathbf{F}\rho h$ (units N\,m$^{-2}$), mapped from the
SnapWave mesh to the two SFINCS cells adjacent to each $u/v$ face, averaged,
rotated into grid coordinates, divided by $\rho_w$ and added as an acceleration
\code{fwuv} in \code{sfincs\_momentum.f90}:
\begin{equation}
\left.\frac{\partial q}{\partial t}\right|_{\mathrm{waves}}
=\phi\;\operatorname{sgn}\!\bigl(\mathrm{fwuv}\bigr)\,
\min\!\left(|\mathrm{fwuv}|,\;F_{\max}\right),
\qquad
F_{\max}=\underbrace{\frac{0.25\sqrt{g}\,\gamma_{\max}^{2}}{T_{p,\mathrm{mean}}}}_{\code{snapwave\_fwmaxfac}/\rho}\,h_{w}^{3/2},
\label{eq:limiter}
\end{equation}
where $\phi$ is the subgrid wet fraction, $h_w$ the local wet depth at the
velocity point, and the cap follows from the depth-limited setup scaling noted
in \code{sfincs\_momentum.f90}
(\code{fwmax = fwmaxfac * hwet * sqrt(hwet)}). Because the default
\code{snapwave\_gammax} is 999, this limiter is effectively \emph{inactive}
unless the user lowers $\gamma_{\max}$; if it is set, the cap grows as
$h_w^{3/2}$ and shrinks with the mean boundary period. The whole force can be scaled by
\code{snapwave\_waveforces\_ratio} (default 1.0; set to 0 to switch off the
incident-wave setup when it is imposed elsewhere, e.g.\ a wavemaker).

The reverse coupling is the depth field: SnapWave recomputes
$h=\max(\zeta_{SFINCS}-z_b,10^{-5})$ at every call, so tide, surge and setup
feed back on wave celerities, breaking and the Herbers IG boundary condition.
Wave-enhanced bed roughness (\code{wave\_enhanced\_roughness}) can additionally
feed $u_{orb}$ back into the SFINCS friction term via the BMI.

\subsection{Stationarity and update interval}

SFINCS updates the wave field every \code{dtwave} seconds (default 3600\,s);
between updates the force field is frozen while the hydrodynamics evolve. The
model therefore assumes that the wave field adapts instantaneously to the water
depth and boundary conditions at each update: there is no wave-memory or
propagation time lag. For domains whose offshore-to-coast travel time
($\sim L/c_g$) is comparable to or larger than \code{dtwave}, or for rapidly
varying storm forcing, this is a structural approximation to be aware of.

\section{Numerical method}
\label{sec:numerics}

Equation~\eqref{eq:energybalance} is solved on an unstructured node mesh with a
locally implicit scheme in $\theta$ and an upwind scheme in space:

\begin{enumerate}
  \item \textbf{Upwind characteristic interpolation.} For each node $k$ and
  direction $\theta_j$, the backward characteristic intersects an edge between
  two surrounding nodes $k_1,k_2$; values are linearly interpolated with weights
  $w_1,w_2$ and distance $\Delta s$ (\code{snapwave\_domain.f90::find\_upwind\_neighbours},
  \code{intersect\_angle}). Spherical geometry is supported via local
  equirectangular scaling.
  \item \textbf{Directional discretisation.} The refraction term uses central
  differences in $\theta$ for interior bins and first-order upwind at the two
  ends of the directional window, giving a tridiagonal system per node:
  \begin{equation}
  A_j=-\frac{c_{\theta,j-1}}{2\Delta\theta},\quad
  B_j=\frac{1}{\dd t^{*}}+\frac{c_g}{\Delta s}+\frac{D_{tot}}{E},\quad
  C_j=\frac{c_{\theta,j+1}}{2\Delta\theta},\quad
  R_j=\frac{\ee^{\mathrm{old}}_j}{\dd t^{*}}
      +\frac{c_g^{\uparrow}\ee^{\uparrow}_j}{\Delta s}
      -S_{ig,j}\,\code{shinc2ig}.
  \label{eq:tridiag}
  \end{equation}
  \item \textbf{Sweeping iteration.} Nodes are visited in four sweep orders
  aligned with $\bar\theta + m\pi/2$ (Gauss--Seidel-like), with the pseudo-time
  term $1/\dd t^{*}$ (\code{snapwave\_dt}, default 36000\,s) acting as
  under-relaxation. Convergence is checked every fourth sweep:
  $\max_k|\ee-\ee^{\mathrm{old}}|/\max \ee <$ \code{snapwave\_crit} (default
  $10^{-3}$), with at most \code{snapwave\_niter} (default 10) four-sweep
  cycles. Fractional dissipation $D_{tot}/E$ is under-relaxed with factor
  \code{snapwave\_relax\_factor\_DoverE} (default 0.25).
  \item \textbf{Sequential splitting.} Within one call the incident balance,
  optional action balance, and IG balance are solved sequentially per node, with
  the IG source computed either once per four sweeps (default,
  \code{snapwave\_iterative\_srcig=0}) or every sweep.
\end{enumerate}

Energy from the previous call is retained as the initial guess (\code{restart}),
which speeds up convergence between slowly varying updates but also means the
solution carries a weak dependence on the previous wave state.


\section{Comparison with Boussinesq and spectral wave models}
\label{sec:comparison}

\subsection{Position in model hierarchy}

\begin{table}[h!]
\centering
\small
\caption{Conceptual comparison of wave-modelling approaches.}
\label{tab:comparison}
\begin{tabular}{@{}p{2.8cm}p{3.9cm}p{4.0cm}p{3.6cm}@{}}
\toprule
 & \textbf{SnapWave (this code)} & \textbf{Spectral (SWAN, XBeach stationary)} & \textbf{Boussinesq / non-hydrostatic} \\
\midrule
Wave field & $\ee(\theta)$ at one $\sigma$ per system (incident + IG) & $N(\sigma,\theta)$, full frequency--direction spectrum & Free surface $\eta(x,y,t)$, phase-resolving \\
Time dependence & Stationary + pseudo-time iteration & Stationary or non-stationary & Fully time-dependent \\
Dispersion & Linear, Eckart approximation & Linear (exact root of $\sigma^2=gk\tanh kh$) & Weakly/fully dispersive (finite-depth terms) \\
Nonlinearity & None in propagation; empirical breaking/IG closures & Quadruplet/triad interactions (optional) & Nonlinear advection and amplitude dispersion \\
Refraction by currents & Not included & Included (action balance) & Intrinsic \\
Diffraction & Not included & Optional phase-decoy / mild-slope variants & Intrinsic (model dependent) \\
Breaking & Baldock et al.\ (1998) & Battjes--Janssen (1978) or saturation models & Emergent (with breaking closures) \\
IG waves & Separate balance, Leijnse et al.\ (2024) source, Herbers BC & Parameterised or not present & Emergent from bound-wave release \\
Wave force & $D/c$-type dissipation closure, Eq.~\eqref{eq:force} & Radiation-stress gradients $-\partial_j S_{ij}$ & Emergent from resolved flow \\
Typical cost & Very low (flood-model compatible) & Moderate & High (event-scale only) \\
\bottomrule
\end{tabular}
\end{table}

\subsection{Key differences from a SWAN-type spectral model}

\begin{enumerate}
  \item \textbf{No frequency dimension.} SnapWave carries one representative
  frequency per wave system; spectral shape is only imposed at the boundary
  (JONSWAP) for the Herbers IG computation. Frequency shifting by currents,
  quadruplet/triad energy transfer and whitecapping are absent. The optional
  action balance changes the \emph{representative} period with wind growth but
  does not create a spectrum.
  \item \textbf{Energy rather than action.} With no ambient currents, energy and
  action balances are equivalent; the difference matters exactly where
  current--wave interaction is important (tidal inlets, river plumes, strong
  alongshore currents). SnapWave will not refract or steepen waves on a current
  field, even though SFINCS computes that current field.
  \item \textbf{Stationarity.} SWAN's non-stationary mode integrates
  $\partial N/\partial t$ physically; SnapWave's $\partial/\partial t^{*}$ is
  only an iteration device. Hysteresis of the wave field (e.g.\ swell arriving
  before local seas) cannot be represented.
  \item \textbf{Wave forcing closure.} SWAN/XBeach feed
  $-\nabla\!\cdot\!\mathbf{S}$ to the flow; SnapWave applies the dissipation
  closure Eq.~\eqref{eq:force}, plus a limiter Eq.~\eqref{eq:limiter}. Setup in
  non-breaking (shoaling) zones and set-down offshore of the breakpoint are
  therefore only crudely represented.
\end{enumerate}

\subsection{Key differences from Boussinesq-type models}

Boussinesq and non-hydrostatic models resolve the phase of the free surface and
retain frequency dispersion through depth-integrated but dispersive momentum
equations. Consequently they generate bound and free infragravity waves, wave
groups, runup and swash \emph{dynamically}, and their momentum forcing is
exactly the resolved wave-averaged flow. SnapWave replaces all of this with
phase-averaged balances and empirical closures (Baldock breaking, Herbers IG
boundary, Leijnse et al.\ IG source). The gain is orders-of-magnitude speed;
the cost is the loss of phase information, wave--wave interaction, harmonic
structure of runup, and reflection (except through the weakly reflective SFINCS
boundary treatment). Note that SFINCS itself also offers a non-hydrostatic
correction (\code{nonh}) and wavemakers; those are alternative, more
phase-resolving routes within the same model, outside SnapWave.


\subsection{Practical implications for coupled simulations}
\label{sec:implications}

Points most relevant when the coupled wave--hydrodynamic results look wrong:

\begin{enumerate}
  \item \textbf{Check where the force comes from.} Because
  Eq.~\eqref{eq:force} is dissipation-driven, wave setup is concentrated where
  $D_w$ (breaking) is large. If your case has little breaking at the boundary
  depth, or breaking is suppressed by $\gamma$/\code{baldock\_ratio} settings,
  the setup will be small regardless of offshore $H_{m0}$.
  \item \textbf{Check the limiter.} The piecewise force cap
  Eq.~\eqref{eq:limiter} is inactive by default ($\gamma_{\max}=999$), but if
  you have lowered $\gamma_{\max}$ it can clip forces strongly for short $T_p$
  or shallow depths; inspect \code{fwmaxfac} if forces look saturated.
  \item \textbf{Check the moving directional window.} The $\theta$-grid is
  re-centred on the energy-weighted mean direction at every call; rapid turning
  of offshore direction between calls can produce abrupt force changes.
  \item \textbf{Check the update interval.} Set \code{dtwave} small enough to
  resolve the surge/tide variation that feeds $h$; recall the wave field is
  frozen between updates.
  \item \textbf{Check depth feedback at boundaries.} The Herbers IG boundary
  uses the \emph{time-varying} depth at the support point
  (\code{find\_nearest\_depth\_for\_boundary\_points}); shallow support points
  ($H_{m0}/h>0.5$) trigger the depth-reset guard and change the IG forcing.
  \item \textbf{Check double counting.} If incident setup is imposed through a
  wavemaker or through the water-level boundary, set
  \code{snapwave\_waveforces\_ratio=0}.
  \item \textbf{Check IG choices.} Compare \code{snapwave\_use\_herbers=1}
  versus 0, \code{snapwave\_ig\_opt=1} versus 2, and
  \code{snapwave\_iterative\_srcig=0} versus 1; the IG source is evaluated from
  the previous sweep/call state, so sensitivities are expected.
  \item \textbf{Check currents.} If wave--current interaction is physically
  important at your site, SnapWave will not represent it; this is a model
  limitation, not a bug.
  \item \textbf{Check convergence.} The log reports \code{error}, \code{\%ok}
  (and IG equivalents) per call; non-converged wave fields
  (\code{ended at iteration ...}) produce inconsistent forces.
\end{enumerate}

\section{Code--equation cross-reference}

\begin{table}[h!]
\centering
\small
\caption{Mapping of equations to source locations.}
\label{tab:map}
\begin{tabular}{@{}lll@{}}
\toprule
Equation / term & Symbol & Source location \\
\midrule
Energy balance \eqref{eq:energybalance} & $\ee$, $E_{\theta,ig}$ & \code{solve\_energy\_balance2Dstat} (\code{snapwave\_solver.f90}) \\
Dispersion \eqref{eq:dispersion} & $k,c_g$ & \code{disper\_approx\_1} (\code{snapwave\_windsource.f90}) \\
Refraction \eqref{eq:ctheta} & $c_\theta$ & \code{compute\_celerities} (idem) \\
Directional spectrum \eqref{eq:cosn} & $m_s$ & \code{update\_boundary\_points} (\code{snapwave\_boundaries.f90}) \\
Breaking \eqref{eq:baldock} & $D_w$ & \code{baldock} (\code{snapwave\_solver.f90}) \\
IG source \eqref{eq:srcig} & $S_{ig}$ & \code{determine\_infragravity\_source\_sink\_term} \\
IG boundary & $H_{m0,ig}$ & \code{compute\_herbers} (\code{snapwave\_infragravity.f90}) \\
Wind input & $S_{wind}$ & \code{windinput} (\code{snapwave\_windsource.f90}) \\
Wave force \eqref{eq:force} & $F$ & final loop of \code{solve\_energy\_balance2Dstat} \\
Force limiter \eqref{eq:limiter} & \code{fwmaxfac} & \code{compute\_snapwave} (\code{sfincs\_snapwave.f90}) \\
Momentum coupling & \code{fwuv} & \code{update\_wave\_field}, \code{sfincs\_momentum.f90} \\
Upwind/tridiagonal \eqref{eq:tridiag} & $A,B,C,R$ & \code{find\_upwind\_neighbours}, \code{solve\_tridiag} \\
\bottomrule
\end{tabular}
\end{table}

\section*{References}
\begin{itemize}
  \item Baldock, T.E., Holmes, P., Bunker, S., van Weert, P. (1998).
  Cross-shore hydrodynamics within an unsaturated surf zone. \emph{Coastal
  Engineering} 34, 173--196.
  \item Battjes, J.A., Janssen, J.P.F.M. (1978). Energy loss and set-up due to
  breaking of random waves. \emph{Proc.\ 16th ICCE}.
  \item Breugem, W.A., Holthuijsen, L.H. (2007). Generalized shallow-water wave
  growth from Lake George. \emph{J.\ Waterway, Port, Coastal, Ocean Eng.}
  \item Herbers, T.H.C., Elgar, S., Guza, R.T. (1994). Infragravity-frequency
  (0.005--0.05 Hz) motions on the shelf. Part I. \emph{J.\ Phys.\ Oceanogr.} 24.
  \item Leijnse, T., van Ormondt, M., Nederhoff, K., van Dongeren, A. (2021).
  Modeling compound flooding in coastal systems using a computationally
  efficient reduced-physics solver. \emph{Coastal Engineering} 163, 103796.
  \item Leijnse, T.\ et al.\ (2024). Estimating nearshore infragravity wave
  conditions at large spatial scales. \emph{Front.\ Mar.\ Sci.} 11:1355095.
  \item Suzuki, T., Zijlema, M., Burger, B., Meijer, M.C., Narayan, S. (2012).
  Wave dissipation by vegetation with layered schematization in SWAN.
  \emph{Coastal Engineering} 59, 64--71.
  \item Booij, N., Ris, R.C., Holthuijsen, L.H. (1999). A third-generation wave
  model for coastal regions (SWAN). \emph{J.\ Geophys.\ Res.} 104(C4).
  \item Eckart, C. (1952). The propagation of gravity waves from deep to shallow
  water. \emph{Gravity Waves}, NBS Circular 521.
\end{itemize}

\end{document}

